\subsection{Vogel projector channels refined by the Kummer-adjacent order-four grading}
\label{subsec:vogel-kummer-z4-refinement}

The exact $E_6$ split-Casimir decomposition
\[
 27\otimes78=1728\oplus351\oplus27
\]
from the Vogel operator certificate is compatible with the adjacent $D_5+U(1)$ grading
inherited from the Kummer order-four grading of $E_8$.  In the $E_6$ Cartan convention used
by the minuscule-$27$ verifier, the $U(1)$ charge functional is
\[
 q=(4,5,6,3,4,2)
\]
in fundamental-weight coordinates, normalized so that
\[
 27=1_4\oplus10_{-2}\oplus16_1,
 \qquad
 78=(45+1)_0\oplus16_{-3}\oplus16_3.
\]
Reducing charges modulo four gives
\[
 27:(1,16,10,0),\qquad 78:(46,16,0,16).
\]
The $351$ summand occurring in $27\otimes78$ is the fundamental module conjugate to
$\wedge^2 27$; its charge profile follows exactly from pairwise sums of the $27^*$ weights:
\[
 \boxed{351:(45,160,130,16)}.
\]
Subtracting this and the $27$ profile from the tensor-product convolution gives
\[
 \boxed{1728:(256,736,576,160)}.
\]
Hence, residue by residue,
\[
 (302,912,716,176)
 =(256,736,576,160)+(45,160,130,16)+(1,16,10,0).
\]
The Vogel spectral projectors are therefore simultaneously resolved by the
Kummer-adjacent order-four charge.  Tensoring by the $A_2$ triplet relevant to the native
CE2 grade multiplies these profiles by three; in particular the $27\otimes3$ channel is
$(3,48,30,0)$, recovering the same $3+48+30$ refinement found directly in the common
$E_8$ grading.

This is a compatibility theorem for concrete $E_6$ representations and projectors.  It
does not identify Kummer geometry with Vogel's conjectural universal Lie algebra or its
diagrammatic $\Lambda$-algebra, and it does not by itself provide the still-missing blind
CE2 sign derivation.
